\section{A creature without a body: substance, presence, and movement}
\label{sec:angel-nature}

\subsection{What remains when bodily parts are removed? (Question 50)}

Imagine removing every bodily feature from a familiar picture of an angel: wings, face, luminous clothing, even the outline that distinguishes a figure from its background. What remains is not an invisible version of the same picture. Aquinas asks us to think of a creature whose acts of understanding and willing belong to no bodily organ. The difficulty is partly our own: imagination presents shapes, whereas intellect can grasp what no shape adequately represents. The inability to picture a spirit is therefore no proof that spirit is impossible. Aquinas begins with the intelligible likeness of creatures to their intelligent Creator, not with an unusually fine kind of matter.\footnote{\ST{I q.~50 a.~1}, corpus and replies; the argument concerns the intelligible perfection of creation.}

An angel is a \emph{substance}: someone who exists, rather than a quality inhering in someone else. But an angel is not God. This immediately raises the hardest objection in article 2. If the angel has no matter to receive a form, is the angel pure actuality, with nothing received or dependent? Aquinas answers by distinguishing two compositions. In a bodily substance, matter and form constitute the nature; the resulting nature also receives existence. An angel lacks the first composition but retains the second. What an angel is does not explain, by itself, that the angel exists. Its essence stands to received existence as potency to act. Here \emph{potency} means capacity for actuality, not merely muscular power or a future possibility.\footnote{\ST{I q.~50 a.~2}, especially ad~3. Compare \emph{De ente et essentia}, cap.~3, beginning \emph{Nunc restat videre}, and cap.~4, beginning \emph{His igitur visis}, in the Baur--Koch text reproduced by Corpus Thomisticum; chapter divisions differ in other editions.}

Aquinas's \emph{De ente et essentia} makes the distinction especially concrete. One can understand what a thing would be while still asking whether it exists. With a created separate substance, immateriality removes matter from its constitution; it does not remove its need for a giver of existence. This is why ``pure spirit'' and ``pure act'' cannot be interchanged. The first may describe an angel; the second belongs to God alone.

Article 3 then breaks the habit of counting angels by their usefulness to us. Aquinas rejects a census limited to the movers needed for the heavens. Spiritual creatures have their own perfection within God's creation. Daniel's immense heavenly court supports their multitude; Aquinas further judges that their number exceeds material multitude, expressing the plenitude of spiritual perfection within creation.\footnote{\ST{I q.~50 a.~3}; Dan.~7:10.}

Why cannot two angels be individuals of one species, as two people are both human? Aquinas answers that, in bodily substances, a common specific form is multiplied in distinct matter. A subsisting immaterial form lacks that principle of multiplication. Thus each angel is specifically distinct. \emph{Species} here names a metaphysical kind, not a biological population; the claim does not make each angel a group or deny personal individuality. John Damascene explicitly leaves their equality or difference in essence to God's knowledge.\footnote{\ST{I q.~50 a.~4}; John Damascene, \emph{Exposition of the Orthodox Faith} II.3.}

Finally, an angel cannot decompose into bodily constituents. Natural incorruptibility means absence of an internal principle of dissolution, not independence from the Creator. A circle drawn in bronze can lose its circular shape because the bronze can take another shape; a subsisting form is not related to an underlying material in that way. God's continuing gift of existence is still indispensable. Aquinas therefore reads Damascene's language of immortality by grace as emphasizing received permanence and, in its fullest sense, immutability in blessedness.\footnote{\ST{I q.~50 a.~5}, especially ad~1 and ad~3; Damascene, \emph{Orthodox Faith} II.3.}

\angeldossier{I q.~50, aa.~1--5}{Defined: God creates the angelic order (Lateran IV, \emph{Firmiter}, DS~800). Church teaching: spiritual nature and immortality (\emph{CCC} 328, 330). Thomist position: no spiritual matter and one angel per species.}{Psalm~103:4 [Hebrew~104:4]; Daniel~7:10; Dionysius, \emph{Divine Names} IV and \emph{Celestial Hierarchy} X, XIV, as invoked by Aquinas; Damascene II.3.}{Aquinas contests Avicebron's material composition in a.~2. Damascene leaves essential equality undecided; Aquinas determines specific diversity.}

\subsection{An appearance at the table (Question 51)}

The visitors to Abraham can be seen and offered food. Must a visible angel therefore have an angelic body? Aquinas distinguishes a body naturally belonging to a creature from a body assumed for a mission. Human bodily life is not a temporary costume: the soul is the body's substantial form, and human knowing begins through the senses. An angel neither animates a body in this way nor needs bodily sensation to complete its natural knowledge. Its assumption of a visible body serves the recipient of the message.\footnote{\ST{I q.~51 aa.~1--2}; Gen.~18--19; Tob.~5 and 12.}

Article 2 argues from the public character of the biblical encounters. Aquinas does not reduce every appearance to an image within one recipient's imagination: people besides the principal recipient see and interact with the visitor. An assumed body is therefore, in the cases under discussion, an external bodily sign. The angel moves it so that it expresses the angel's spiritual powers. This is a union for representation and action; it is not the union by which a human soul constitutes a living human body. Nor is it the personal union of the divine Word with human nature in the Incarnation. Aquinas regards the earlier appearances as ordered towards that unique coming of the Son.\footnote{\ST{I q.~51 a.~2}, corpus and ad~1--2.}

Article 3 tests the distinction against walking, speaking, seeing, eating, and generation. Producing articulate sound does not require that the producer have a living human voice. Moving legs does not make the assumed body an organism. Its eyes signify intellectual perception; the angel does not acquire knowledge through them. Apparent eating is not nourishment of an angelic metabolism. Raphael's disclosure in Tobias 12:19 supplies Aquinas's immediate biblical comparison. The resurrection of Christ is different: Christ possesses a true living human body, not an angel's representational instrument.\footnote{\ST{I q.~51 a.~3}, especially ad~2--5.}

These distinctions also govern the article's treatment of Genesis 6. Aquinas follows the reading of the ``sons of God'' as Seth's descendants and denies that angels generate offspring by an angelic bodily nature. His accompanying physical explanation of certain inherited reports depends on medieval natural philosophy. Article 2 likewise proposes condensed air as the material of an assumed body.\footnote{\ST{I q.~51 a.~2 ad~3; a.~3 ad~6}.}

\angeldossier{I q.~51, aa.~1--3}{Church teaching: angels are purely spiritual (\emph{CCC} 330). Thomist position: the analysis of assumed bodies and their operations. Disputed: the proposed physical constitution of such bodies.}{Dionysius, \emph{Divine Names} IV; Augustine, \emph{City of God} XVI, cited in a.~2; Genesis~18--19; Tobias~12:19.}{Aquinas rejects reducing all appearances to inward visions. The bodily accounts attributed to Origen and other early witnesses in a.~1 require their own historical distinctions; they are not identical positions.}

\subsection{Presence without size (Question 52)}

A person occupies a room because a body has dimensions. If an angel has no dimensions, can the angel be there at all? Aquinas begins with the Church's prayer for the angels' presence, then asks what sort of presence makes sense. A body contacts a place through its measurable extension. An angel is present through the application of its power. The angel is not squeezed into a spiritual point: a geometrical point belongs to quantity, whereas an angel does not.\footnote{\ST{I q.~52 aa.~1--2}. The \emph{sed contra} of a.~1 invokes the collect beginning \emph{Visita, quaesumus}; its appearance here reports Aquinas's authority, not an edition-specific liturgical prescription.}

The traditional terms help if their definitions remain attached. A body is present \emph{circumscriptively}: its dimensions correspond to the place. An angel is present \emph{definitively}: its operation determines a here, without bodily measurement, and excludes unrestricted presence elsewhere. God is present everywhere by his universal causality. The creature's lack of size never becomes the Creator's omnipresence. An angel's one field of immediate operation can extend over something divisible; ``one place'' is not necessarily the footprint of a human body.\footnote{\ST{I q.~52 a.~2}, corpus and replies.}

Article 3 is consequently subtler than the familiar question of how many spirits fit into a tiny space. Aquinas denies two angels the same place because, on his account, each would have to be the complete immediate operative cause in the same respect. He does not argue that they collide. Several rowers can move one boat because each contributes partial power; they act together as one sufficient mover. Nor are a soul and an angel related to a body in the same way: the soul is its form. His conclusion follows from his causal definition of angelic location.\footnote{\ST{I q.~52 a.~3}, corpus and ad~1--3.}

\angeldossiersettled{I q.~52, aa.~1--3}{Thomist position: presence by application of finite power, and the causal account of exclusive location.}{The collect \emph{Visita, quaesumus}, invoked in a.~1; Damascene, \emph{Orthodox Faith} II.3, for finite presence; the soul--body analogy in a.~3.}

\subsection{A journey with a before and an after (Question 53)}

If an angel is present by acting here, movement means a succession of such contacts. It need not mean that invisible material travels through the air. Aquinas allows a continuous passage, in which the operative contact withdraws successively from one region and reaches another, and a discontinuous change from one location to another. In the second case passage through every intervening place is unnecessary. The distinction concerns what an immaterial agent's presence requires; it supplies no angelic speed in miles per hour.\footnote{\ST{I q.~53 aa.~1--2}.}

Neither possibility abolishes succession. Article 3 rejects the claim that the entire movement occurs in one instant. There must be a before and an after: the angel is first here and then there. Aquinas calls the measure of that succession time, while distinguishing it from the continuous time measured by heavenly motion in his cosmology. A discontinuous succession of angelic acts is not simply a shorter interval on the same physical clock.\footnote{\ST{I q.~53 a.~3}, especially ad~1 and ad~3.}

This also answers a spiritual puzzle: why would a blessed angel move if happiness leaves it nothing to acquire? Aquinas replies that ministry answers another's need. Movement need not arise from an agent's misery or imperfection; a perfected creature can exercise its fullness for someone else. Its journey does not interrupt the vision of God.\footnote{\ST{I q.~53 a.~1 ad~2--3}; Heb.~1:14.}

\angeldossier{I q.~53, aa.~1--3}{Thomist position: continuous or discontinuous movement, always with succession.}{The descent of Christ's soul, invoked by analogy in a.~1; Aristotle's account of movement and time; Hebrews~1:14.}{Aquinas rejects instantaneous movement understood as eliminating succession, while admitting a discontinuous transition without intervening spatial contacts.}
